% !TeX root = main.tex
\documentclass[main.tex]{subfiles}

\begin{document}

\noindent
Malware is one of the most significant threats faced by users in the current digital landscape. The increasingly sophisticated nature of these programs, in terms of both their capabilities and their stealth, mean that this threat is only increasing as we become evermore reliant on digital infrastructure in our daily lives. 

Therefore, detecting malware before it can cause damage is of high priority to security professionals. However, following the introduction of anti-virus software, malware developers have begun using various evasion techniques to avoid detection. This has led to a cat-and-mouse like game, where cyber-criminals adopt new forms obfuscation, and security researchers develop new techniques to handle them. This makes conventional approaches to malware detection challenging, due to its constantly evolving nature. 

Machine learning however offers a potential means of dealing with this by finding ways to detect malware in spite of any evasion techniques. In this investigation, we evaluate various approaches across a number of machine learning domains, such as linear modelling, natural language processing, and image processing to investigate both their performance at detecting windows malware, and their reliance on obfuscation. 

We develop a RoBERTa-based language model trained on sequences of assembly instruction, and a Convolutional Neural Network using image representations of programs to identify malware. These are both considered static approaches whereby features are extracted from the structure of the program, without executing it. Dynamic methods on the other hand, observe the behaviour of a program, and are typically more capable than static methods. However, they often incur a significant overhead in their complexity, as an environment where potentially malicious software can be safely executed is needed. This is typically a remote server hosting a virtual machine that isolates the programs execution, preventing it from causing harm. However, to communicate with this sandbox, an internet connection is required. In comparison, static approaches are able to operate in environments where an internet connection is not guaranteed. Therefore, in light of their disadvantages, static methods are still relevant in domain. 

Due to the highly pervasive nature of obfuscation in modern malware, we theorise that static approaches rely heavily on the presence of obfuscation as an indicator for malware. However, obfuscation is not exclusive to malicious programs, as legitimate software developers may use it to protect their intellectual property from reverse engineering. Therefore, to better understand how static approaches identify malware, we conduct an additional experiment that aims to evaluate how our models distinguish between benign, and malicious programs. 

The research questions we seek to answer in this work are as follows: 

\begin{enumerate}[label=\textbf{RQ\arabic*}:]
    \item \emph{Are modern natural language processing techniques, such as language models and pretraining effective at identifying malware?}
    \item \emph{Does pretraining improve the accuracy and reliability of malware detection models based on techniques from natural language processing?}
    \item \emph{To what extent are static malware detection models resilient to obfuscation?}
\end{enumerate}


\section{Professional Considerations}

This research is conducted in strict accordance with the BCS Code of Conduct~\citep{britishcomputingsocietyBCSCodeConduct2022}, ensuring adherence to the ethical, legal, and professional standards expected of computing professionals. The work is structured to protect the public interest, maintain professional competence and integrity, uphold obligations to the relevant authority, and safeguard the reputation of the profession.

\subsubsection{Public Interest}

We conduct our research with regard to the safety, rights, and welfare of the public to ensure our work is in their interest. Samples of malicious programs will be handled in a secure environment, with safeguards to protect against execution on both our own hardware, and that of others to avoid accidental infection. Furthermore, proprietary programs subject to copyright law will not be distributed to ensure the intellectual property rights of third parties are respected. Our research will be conducted without discrimination on the basis of ethnicity, sex, religion, age or other protected characteristics, and we distribute all materials used to conduct this investigation to support transparency, reproducibility, and to benefit the wider research community. 

To ensure the malware used in this investigation are unable to infect any machines used to conduct our research, we prevent them from being executed using the \texttt{chmod} program. This is a tool in UNIX based systems that modifies the permissions associated with a file, and can prevent files from being executed on the system. 

\subsubsection{Professional Competence and Integrity}


We conduct our research within the verified level of competence of the authors, and avoid misrepresenting our skills and expertise. The project itself contributes to the continued professional development of our skills in Artificial Intelligence, and Cybersecurity. We employ established best practices in both fields where appropriate, and comply with all relevant UK legislation, including data protection, intellectual property, and cyber-security. Diverse perspectives are actively sought, and the authors welcome constructive criticism to ensure the quality and robustness of the work. No element of this research will cause harm to the property, reputation, or the livelihood of others through false or negligent actions.

\subsubsection{Duty to the Relevant Authority}


To ensure the duty of the authors to relevant authorities, we conduct this research in compliance with the policies and procedures of the University of Sussex, and any other governing bodies. Any potential conflicts of interest are avoided, and disclosed in cases where they are unavoidable. We accept any professional responsibilities for our work, and that of others involved with this project. Results and findings will be reported accurately, and truthfully, without distortion, omission, or misrepresentation. 

\subsubsection{Duty to the Profession}
 

The authors acknowledge their responsibilities to uphold the credibility and reputation of the computing profession, and we pledge to take no actions that will take it into disrepute.

\end{document}